% BHU PYQ -- Manually transcribed & error-corrected
% Paper: MTB-402 : Mathematical Methods
% Exam: B.Sc. (Hons) Semester IV Examination, 2013-14

\documentclass[11pt,a4paper]{article}
\usepackage[a4paper,top=2.5cm,bottom=2.5cm,left=2.8cm,right=2.8cm,headheight=14pt]{geometry}
\usepackage{amsmath,amssymb}
\usepackage[T1]{fontenc}
\usepackage[utf8]{inputenc}
\usepackage{lmodern}
\usepackage{microtype}
\usepackage{fancyhdr}
\usepackage{enumitem}
\usepackage{hyperref}
\usepackage{lastpage}
\usepackage{parskip}

\hypersetup{colorlinks=true,urlcolor=black,linkcolor=black,
  pdftitle={MTB-402: Mathematical Methods -- BHU B.Sc. Sem IV 2013-14}}
\pagestyle{fancy}\fancyhf{}
\renewcommand{\headrulewidth}{0.4pt}\renewcommand{\footrulewidth}{0.4pt}
\fancyhead[L]{\small   \textit{Banaras Hindu University}}
\fancyhead[R]{\small   \textit{MTB-402: Mathematical Methods}}
\fancyfoot[C]{\small Page  \thepage\ of \pageref{LastPage}}


\newlist{parts}{enumerate}{2}
\setlist[parts,1]{label=(\alph*),leftmargin=*,itemsep=5pt,topsep=3pt}
\setlist[parts,2]{label=(\roman*),leftmargin=*,itemsep=3pt,topsep=2pt}

\newcommand{\pts}[1]{\hfill   \textit{[#1]}}


\begin{document}

\begin{center}
  {\large\bfseries BANARAS HINDU UNIVERSITY}\\[2pt]
  {
\normalsize B.Sc. (Hons) --- Semester IV Examination, 2013-14}\\[6pt]
  \rule{0.8\linewidth}{0.5pt}\\[6pt]
  {\large\bfseries Paper No. MTB-402: Mathematical Methods}\\[4pt]
  {
\normalsize Time: 3 Hours \hfill Full Marks: 70}
\end{center}
\rule{\linewidth}{0.4pt}
\smallskip



\noindent   \textit{Instructions: Answer   \textbf{five} questions including Question~1 which is compulsory. Figures in right-hand margin indicate marks.}
\bigskip




\noindent   \textbf{Question 1.}   \textit{(Compulsory.)}

\begin{parts}
  \item Using the definition of the Laplace transform, find $\mathcal{L}\{\sin at\}$. \pts{2}
  \item Prove that the Laplace transformation is linear. \pts{2}
  \item Using the convolution theorem, find the value of $\mathcal{L}^{-1}\!\left\{\dfrac{1}{s(s+a)}\right\}$. \pts{3}
  \item Find the Laplace transform of $f(t)$, where $f(t) = e^t$ for $0 \le t < 2\pi$ and $f(t + 2\pi) = f(t)$. \pts{3}
  \item Write the function whose Laplace transform is $1$. \pts{1}
  \item If $\mathcal{L}^{-1}\{F(s)\} = f(t)$, then prove that $\mathcal{L}^{-1}\!\left\{F(as)\right\} = \dfrac{1}{a}f\!\left(\dfrac{t}{a}\right)$. \pts{2}
  \item If $f(x)$ has the Fourier coefficients $a_n, b_n$, show that $kf(x)$, where $k$ is a constant, has the Fourier coefficients $ka_n, kb_n$. \pts{2}
  \item If the Fourier cosine integral of $f(x)$ is given by $f(x) = \int_0^\infty A(\omega)\cos(\omega x)\,d\omega$, prove that $-\dfrac{d^2}{d x^2}f(x) = \int_0^\infty \omega^2 A(\omega)\cos(\omega x)\,d\omega$, where $A(\omega) = \dfrac{2}{\pi}\int_0^\infty f(x)\cos(\omega x)\,dx$. \pts{2}
  \item Define the $n$-th order proximity of two curves. \pts{2}
  \item Test for an extremum for the functional $I[y] = \int_0^1 (y'^2 + y^2 - 2y'^2 y')\,dx$, $y(0) = 1$, $y'(0) = 2$. \pts{3}
\end{parts}

\medskip\rule{\linewidth}{0.2pt}

\medskip


\noindent   \textbf{Question 2.}

\begin{parts}
  \item If $f(t)$ is a continuous function except for an ordinary discontinuity at $t = a$ ($a > 0$), then show that:
  \[
  \mathcal{L}\!\left\{\frac{d}{dt}f(t)\right\} = s\,\mathcal{L}\{f(t)\} - f(0^+)
  \] \pts{6}
  \item Express $f(t)$ in terms of unit step functions and hence find its Laplace transform, where:
  \[
  f(t) = 
\begin{cases} 0, & 0 < t < 1 \\ t-1, & 1 < t < 2 \\ 1, & t > 2 \end{cases}
  \] \pts{6}
\end{parts}

\medskip\rule{\linewidth}{0.2pt}

\medskip


\noindent   \textbf{Question 3.}

\begin{parts}
  \item State and prove the existence theorem for the Laplace transform. \pts{6}
  \item Using Laplace transforms, solve the differential equation:
  \[
  \frac{d^2y}{dx^2} + 2\frac{dy}{dx} - y = x, \quad y(0) = 0,\quad y'(0) = 1
  \] \pts{6}
\end{parts}

\medskip\rule{\linewidth}{0.2pt}

\medskip


\noindent   \textbf{Question 4.}

\begin{parts}
  \item Find $\mathcal{L}^{-1}\!\left\{\dfrac{\ln(s+a)}{s+b}\right\}$. \pts{6}
  \item Find the Fourier cosine series expansion of the periodic function of period $2P$, defined by $f(t) = \sin\!\left(\dfrac{\pi t}{P}\right)$ for $0 < t < P$. \pts{6}
\end{parts}

\medskip\rule{\linewidth}{0.2pt}

\medskip


\noindent   \textbf{Question 5.}

\begin{parts}
  \item Using Laplace transforms, solve the integral equation:
  \[
  y(t) = \sin 2t + \int_0^t y(x)\sin 2(t-x)\,dx
  \] \pts{6}
  \item Find the Fourier integral representation of the function:
  \[
  f(x) = 
\begin{cases} 0, & x < 0 \\ 1, & 0 < x < 1 \\ 0, & x > 1 \end{cases}
  \]
  Hence find the value of $\int_0^\infty \dfrac{\sin\omega}{\omega}\,d\omega$. \pts{6}
\end{parts}

\medskip\rule{\linewidth}{0.2pt}

\medskip


\noindent   \textbf{Question 6.}

\begin{parts}
  \item Find $\mathcal{L}^{-1}\!\left\{\dfrac{5s+3}{(s-1)(s^2+2s+5)}\right\}$. \pts{6}
  \item Using the calculus of variations, find the shortest distance between the point $(1, 0)$ and the ellipse $4x^2 + 9y^2 = 36$. \pts{6}
\end{parts}

\medskip\rule{\linewidth}{0.2pt}

\medskip


\noindent   \textbf{Question 7.}

\begin{parts}
  \item Find the extremal for which the functional $I = \int_0^1 y'^2\,dx$, $y(0) = 0$, $y(1) = 1$ is an extremum, subject to the isoperimetric condition $\int_0^1 y\,dx = \dfrac{1}{6}$. \pts{6}
  \item Find the extremal for the functional:
  \[
  I = \int_0^{\pi/2} \left(y''^2 - y^2 + x^2\right)dx, \quad y(0) = 0,\quad y'\!\left(0\right) = 1,\quad y\!\left( frac{\pi}{2}\right) = y'\!\left( frac{\pi}{2}\right) =  frac{1}{4}
  \] \pts{6}
\end{parts}

\vfill
\rule{\linewidth}{0.4pt}

\begin{center}\small
  \href{https://bhu-exam.vercel.app/}{Click here to visit our website}\\[4pt]
  \href{https://me-aryan.vercel.app/}{Click here to visit our contributor}\\[4pt]
  \href{https://quashan-me.vercel.app/}{Click here to visit our contributor}
\end{center}

\end{document}
